\section*{الفصل \ref{ch:g8:the-oceans} --- المحيطات}

\begin{solution}{exo:g8:the-oceans:1}
\omterm{def:g8:the-oceans:seafloor}{الرف القاري}، ثم \omterm{def:g8:the-oceans:seafloor}{المنحدر القاري}، ثم \omterm{def:g8:the-oceans:seafloor}{السهل السحيق}، ثم \omterm{def:g8:the-oceans:seafloor}{الظهرة المحيطية}.
\end{solution}

\begin{solution}{exo:g8:the-oceans:2}
(أ) خطأ: نحو \qty{3700}{m}، أي \qty{3.7}{km}. (ب) صواب. (ج)
خطأ: التبخر يزيل الماء ويترك الأملاح، فيصير الماء أكثر
\omterm{def:g8:the-oceans:salinity}{ملوحة}. (د) صواب.
\end{solution}

\begin{solution}{exo:g8:the-oceans:3}
(أ) \omterm{def:g8:the-oceans:seafloor}{ظهرة محيطية}. (ب) \omterm{def:g8:the-oceans:salinity}{الملوحة}. (ج) \omterm{def:g8:the-oceans:current}{دوامة محيطية}. (د) \omterm{def:g8:the-oceans:seafloor}{خندق محيطي}.
\end{solution}

\begin{solution}{exo:g8:the-oceans:4}
(أ) عكس اتجاه عقارب الساعة. (ب) الدافئ: تيار الخليج؛ والبارد: تيار الكناري.
(ج) \omterm{def:g8:winds-global-circulation:global-winds}{الرياح التجارية}.
\end{solution}

\begin{solution}{exo:g8:the-oceans:5}
(أ) $35 \div 1000 \times 100 = 3.5~\%$. (ب) $20 \times 35 = \qty{700}{g}$.
(ج) $1000 \div 35 \approx 28.6$: أي نحو \qty{29}{kg}.
\end{solution}

\begin{solution}{exo:g8:the-oceans:6}
(أ) $1500 \times 6 = \qty{9000}{m}$. (ب) يهبط الصوت إلى القاع ثم يعود
صاعدًا: فهو يقطع العمق مرتين. (ج) $9000 \div 2 =
\qty{4500}{m}$. (د) \omterm{def:g8:the-oceans:seafloor}{سهل سحيق} (فالخندق أعمق، والظهرة أو الرف
أقل عمقًا).
\end{solution}

\begin{solution}{exo:g8:the-oceans:7}
(أ) بين نحو \qty{3500}{m} و\qty{5000}{m}. (ب) على بعد نحو \qty{4400}{km} على طول المقطع، وعلى عمق نحو \qty{2300}{m}. (ج) نحو
$4500 - 2300 = \qty{2200}{m}$ (بين 1200 و\qty{2700}{m}، بحسب \omterm{def:g1:landscapes-around-us:landform}{السهل}
المختار). (د) عمق \omterm{def:g2:water-on-earth:ocean}{المحيط} بضعة كيلومترات لكن عرضه آلاف الكيلومترات: فلو رُسم
بالمقياس نفسه، لكان المقطع كله أرق من خط.
\end{solution}

\begin{solution}{exo:g8:the-oceans:8}
\omterm{def:g2:water-on-earth:ocean}{البحر} A أكثر \omterm{def:g8:the-oceans:salinity}{ملوحة}: فالشمس والهواء الجاف يسببان \emph{تبخرًا} قويًا،
فيرحل الماء و\emph{تبقى الأملاح}، ولا يكاد \omterm{def:g4:clouds-rain-snow:precipitation}{مطر} أو \omterm{def:g1:landscapes-around-us:water}{نهر} يعيد
\emph{\omterm{def:g2:water-on-earth:salt-fresh}{الماء العذب}}. \omterm{def:g2:water-on-earth:ocean}{والبحر} B أقل \omterm{def:g8:the-oceans:salinity}{ملوحة}: فأنهاره ومطره تجلب كثيرًا من
\emph{\omterm{def:g2:water-on-earth:salt-fresh}{الماء العذب}}، الذي يخفف تركيز الأملاح.
\end{solution}

\begin{solution}{exo:g8:the-oceans:9}
(أ) تيار الخليج، نحو الشمال على طول \omterm{def:g1:landscapes-around-us:water}{الساحل}، ثم نحو الشمال الشرقي. (ب) نحو
الشرق (الشمال الشرقي)، مع تيار شمال الأطلسي. (ج) على طول \omterm{def:g1:landscapes-around-us:water}{ساحل} شمال غرب
إفريقيا (وإسبانيا والبرتغال)، مع تيار الكناري. (د)
بارد.
\end{solution}

\begin{solution}{exo:g8:the-oceans:10}
(أ) $6.4 \times 24 = 153.6$: أي نحو \qty{154}{km}. (ب) $153.6 \times 7
\approx \qty{1075}{km}$. (ج) تيار الخليج أسرع قليلًا من الماشي، لكنه لا
يتوقف أبدًا، ليلًا ونهارًا.
\end{solution}

\begin{solution}{exo:g8:the-oceans:11}
(أ) \qty{3.7}{km} و\qty{11}{km}. (ب) $11000 \div 3700 \approx 3.0$
مرات. (ج) يقطع الصوت $2 \times 11000 = \qty{22000}{m}$؛ $22000
\div 1500 \approx 14.7$: أي نحو \qty{15}{s}. (د) $1000 = 10^3$ سنة.
\end{solution}

\begin{solution}{exo:g8:the-oceans:12}
(أ) قرب \omterm{def:g3:continents-and-oceans:equator}{خط الاستواء}، تدفع \omterm{def:g8:winds-global-circulation:global-winds}{الرياح التجارية} الماء المداري الدافئ نحو الغرب.
وعند الحافة الغربية للأطلسي ينعطف شمالًا مع تيار الخليج، وهو تيار دافئ، على
طول فلوريدا وأمريكا الشمالية. ثم تدفعه \omterm{def:g8:winds-global-circulation:global-winds}{الرياح الغربية} عبر الأطلسي نحو الشمال
الشرقي، حتى البحار القريبة من النرويج. (ب) هناك يعطي حرارته للهواء البارد:
فيبرد، ويزداد \omterm{def:g8:the-oceans:salinity}{ملوحة} حين يتكون الجليد البحري، لأن الملح يبقى في الماء.
والماء البارد المالح كثيف، فيغوص إلى القاع. (ج) يزحف جنوبًا على طول قاع
\omterm{def:g2:water-on-earth:ocean}{المحيط}، نحو القارة القطبية الجنوبية، ثم إلى المحيطين الهندي والهادئ، حيث
يصعد ببطء من جديد.
\end{solution}

\begin{solution}{pb:g8:the-oceans:1}
\textbf{1.} في \omterm{def:g2:water-on-earth:ocean}{المحيط} الأطلسي، في \omterm{def:g3:continents-and-oceans:equator}{نصف الكرة الشمالي}.

\textbf{2.} مع اتجاه عقارب الساعة؛ وتُحمل العوّامة شمالًا في البداية.

\textbf{3.} ماء المرحلة A (تيار الخليج) دافئ: فهو آتٍ من المناطق المدارية.
وماء المرحلة C (تيار الكناري) بارد: فهو آتٍ من عروض أعلى، متجهًا نحو \omterm{def:g3:continents-and-oceans:equator}{خط
الاستواء}.

\textbf{4.} المرحلة B: \omterm{def:g8:winds-global-circulation:global-winds}{الرياح الغربية}. والمرحلة D: \omterm{def:g8:winds-global-circulation:global-winds}{الرياح التجارية}.

\textbf{5.} $3800 \div 100 = 38$ يومًا.

\textbf{6.} B: $2100 \div 20 = 105$ أيام. C: $3100 \div 20 = 155$ يومًا.
D: $6100 \div 20 = 305$ أيام.

\textbf{7.} $3800 + 2100 + 3100 + 6100 = \qty{15100}{km}$.

\textbf{8.} تيار الخليج هو التيار الضيق السريع على الجانب الغربي من \omterm{def:g8:the-oceans:current}{الدوامة}؛
أما التيارات على جوانبها الأخرى فعريضة بطيئة.

\textbf{9.} $38 + 105 + 155 + 305 = 603$ أيام.

\textbf{10.} $603 \div 30 = 20.1$: أي نحو \textbf{20 شهرًا}.

\textbf{11.} بعد المراحل A وB وC، يكون قد مضى $38 + 105 + 155 = 298$ يومًا؛
ويحل اليوم 365 أثناء المرحلة D (التيار الاستوائي الشمالي)، بعد $365 -
298 = 67$ يومًا من بدايتها.

\textbf{12.} $603 - 365 = 238$ يومًا إضافيًا، أي $238 \div 30 \approx 8$
أشهر إضافية.
\end{solution}
